\scanpage{141}
La catégorie des groupes finis \inserted{comm. $G$} sur $K$ tel que pour tt $\ell$, $G(\ell)$ soit non ramifié rel. à $S[1/\ell]$, est une sous catégorie épaisse de celle des motifs sur $S$, la catégorie quotient est équivalente à celle des isomotifs lisses sur $S$.

La catégorie des motifs est une $\otimes$-catégorie avec \uline{Hom}. Les foncteurs motifs $\longrightarrow$ isomotifs, motifs $\to$ faisceaux $\ell$-adiques, sont des $\otimes$-foncteurs, compatibles aux \uline{Hom}.

Un \deleted{définit} hom de motifs est une isogénie s'il induit un isom sur les isomotifs sousjacents. On définit la dimension (ou le poids), le niveau etc en termes des isomotifs.

Tout schéma propre et lisse $X$ sur $S$ définit des motifs\inserted{ sur $S$}, notés $R^i f_*(\mathbb Z_X)$ s'il n'y a pas de confusion à craindre. de $=$ pour des \unsure{sites} étales relatifs lisses et propres.

\uline{Relations avec les schémas en groupes commutatifs propres et lisses}.

On définit un foncteur \uline{contravariant} de la catégorie de ces objets, soit $A$, dans la catégorie des motifs sur $S$, en posant
\[
\begin{aligned}
&\deleted{\unsure{\Phi|\varphi|\mathfrak P}(A)=R^1 f_*(\mathbb Z_A)=\bigl(\ R^1 f_*(\bigr.}\\
&\Phi(A)=(M,E_\ell,u)
\end{aligned}
\]
où
\[
\left\{\begin{aligned}
 M&=R^1(f_{A^0})_*(\mathbb Q_{A^0})\\
 E_\ell&=\underline{\operatorname{Hom}}({}_{\ell^\infty}A,\mathbb Q_\ell/\mathbb Z_\ell)\\
 &u_\ell\ \text{provenant de l'hom évident}\\
 &\hspace{4em}E_\ell\longrightarrow R^1 f_{A^0*}(\mathbb Z_{\ell A^0})
\end{aligned}\right.
\]
\margin{extension de $\underline{\operatorname{Hom}}(A/A^0,\mathbb Q_\ell/\mathbb Z_\ell\unsure{)}$ par $R^1 f_{A^0*}(\mathbb Z_{\ell A^0})$}

\uline{NB} Il suffirait que $A$ \deleted{\unsure{soit|soit un} \unsure{à supp.|extension}} \inserted{soit extension d'un schéma} en gr fini $A/A^0$ par un schéma abélien $A^0$, avec $A/A^0(\ell)$ extension par $0$ d'un groupe fini étale sur $S_\ell\deleted{\unsure{0}}$.
% archive sheet number: 140
